\section{Limitations}
\label{sec:limitations}

Our evidence comes from six MOT-Challenge training-split videos
\cite{leal2015mot15,milan2016mot16}; there is no held-out test-split HOTA,
since the public benchmarks withhold test ground truth and we required it for
the per-clip detector-precision and association-isolation analyses. We
therefore report our findings as \textbf{per-clip and correlational rather
than as a universal law}: the link between low detector precision on dense,
low-resolution clips and the ``more detections hurt'' regime is consistent
across the sequences we audit (\cref{tab:detector,fig:money}), but six clips
cannot establish a calibrated precision threshold that transfers to arbitrary
domains. The practical guideline of \cref{sec:tuning} should be read as a
tuning heuristic, not a tuned constant.

A reviewer may object that restoring the original DeepSORT confidence and NMS
gate is ``self-inflicted'': we re-broke a filter the method already had, then
claimed credit for re-adding it. We address this head-on. The contribution is
not the gate itself but the \emph{audit} of what a routine modernization
silently discards: swapping in a modern detector and re-identification model
\cite{jocher2023yolov8,zhou2019osnet} quietly drops the gating contract, and
the resulting regression is invisible in mean HOTA yet large and predictable
per clip. The finding is also specific in mechanism. It concerns ungated
low-confidence detections \emph{instantiating spurious new tracks}, which is
distinct from ByteTrack \cite{zhang2022bytetrack}, where low-confidence boxes
are retained but only \emph{associated} to existing tracks and never seed new
ones. That distinction is what makes the result a property of the gating
contract rather than a tautology.

Detector and re-identification coverage is narrow. Two additional detector
backends, NanoDet and RTMDet, were unbuildable on our runtime (the NumPy-2.x
and ground-truth column-format issues noted in \cref{sec:method}), so the
detection-side claims rest on a single live detector and cannot speak to
cross-detector generality. Likewise the live pipeline uses one
re-identification family; the OSNet-versus-timm comparison
\cite{zhou2019osnet,wightman2019timm} isolates association only under
ground-truth boxes, so appearance robustness under real detections, and across
other re-identification architectures, remains untested. Confirming the
precision-gating guideline against modern DeepSORT descendants
\cite{du2023strongsort,cao2023ocsort,aharon2022botsort} on full benchmark
splits is left to future work.
